


% \textcolor{red}{Cas Notes on key TODOS:}
% \begin{itemize}
%     \item \textcolor{red}{Clear paper with AISI. Figure out framing and redactions. Get green light from authors.}
%     \item \textcolor{red}{Make sure we are ready to release the models, training code, attack code. Probably don't release modernbert, blocklist, and filtered docs. Mention somewhere about how the Weak/Strong filter models with CB+LAT are the two most robustly bio-ignorant models overall.}
%     \item \textcolor{red}{Table of models. Cref it.}  
%     \item \textcolor{red}{Comb through.} 
%     \item \textcolor{red}{Final fig sizing and spacing. Check for errors and bad cites.}
% \end{itemize}

% \textcolor{darkgreen}{Kyle Notes on key TODOS:}
% \begin{itemize}
%     \item \textcolor{darkgreen}{Resolve comments in Section 2}
%     \item \textcolor{darkgreen}{Resolve Stella's feedback}
% \end{itemize}

\newpage
\section{Introduction} \label{sec:intro}

% Topic: Models are getting a lot better at dangerous capabilities
As frontier large language models (LLMs) grow more advanced, their developers have raised concerns about increasing security risks posed by their models. 
For example, in its Gemini 2.5 Pro Preview model card, Google Deepmind remarks, ``\textit{subsequent revisions in the next few months could lead to a model that reaches the critical capability level [for harmful novice uplift]}'' \citep{google2025gemini2_5_pro_preview}.
Anthropic reported that it was preemptively activating its \textit{Safety Level 3} protocols for Claude Opus 4, writing that they ``\textit{could not rule out}'' ``\textit{the ability to significantly assist individuals or groups with basic STEM backgrounds in obtaining, producing, or deploying CBRN\footnote{Chemical, Biological, Radiological, and Nuclear (CBRN)} weapons.}'' \citep{anthropic_ASL3_2025}. 
% These examples illustrate the pressing need for robust safeguards against misuse.
% OpenAI's Deep Research system card reported, ``\textit{several of our biology evaluations indicate our models are on the cusp of being able to meaningfully help novices create known biological threats.}'' \citep{OpenAI2025}. 
OpenAI's ChatGPT Agent system card states ``\textit{We have decided to [precautionarily] treat this launch as high capability in the Biological and Chemical domain}'' \citep{OpenAI2025_ChatGPT_Agent}.

% \stella{I would rather frame this as ``current approaches to AI safety are overfit to the deployment context of closed-weights models and don't consider the context and needs of the open ecosystem.''} Frontier AI systems, such as those mentioned above, are often deployed with closed weights. 
% Open-weights models, however, are being released and modified at an increasing rate \citep{bhandari2025forecasting}, and their capabilities tend to lag only six to twelve months behind those of closed-weights models \citep{epoch2024openmodelsreport, maslej2024artificial}.
% Open-weights models offer unique benefits related to transparency, open research, and the de-concentration of power \citep{bommasani2024considerations, kapoor2024societal, françois2025differentapproachaisafety}.
% However, they also create unique risks from downstream modifications and unmonitored use \citep{seger2023open, chan2023hazards, huang2024harmful}.
% For these reasons, managing risks from open-weights models is a key challenge toward realizing the benefits and mitigating the harms of AI overall \citep{Bengio2025InternationalAS}.

Today, frontier LLMs such as the above are often deployed with closed weights behind APIs.
However, open-weight models are being released at an increasing rate \citep{bhandari2025forecasting}, and their capabilities tend to lag only six to twelve months behind those of closed-weight models \citep{epoch2024openmodelsreport, maslej2024artificial}.
Open models offer unique benefits related to transparency, research, and the deconcentration of power \citep{bommasani2024considerations, kapoor2024societal, eiras2024risks, françois2025differentapproachaisafety}. 
However, they create unique risks from downstream modifications and unmonitored use \citep{seger2023open, chan2023hazards, huang2024harmful, eiras2024risks}. 
This prompts the question: How can harmful uses of advanced open-weight models be effectively mitigated?

Despite the rising prominence of open-weight models, the science of open-weight model safety is nascent \citep{françois2025differentapproachaisafety, pai2024risk}. 
Closed-weight deployments keep a model's weights secure and allow for monitors and filters on its inputs and outputs.
However, \textbf{the defining challenge for open-weight model safety is that open models can be modified arbitrarily by downstream actors}.
% As a result, open-weight model risk management is challenging, and progress has not been commensurate with the increasing risks and opportunities from open-weights models. 
Despite recent work on developing \textit{tamper-resistant} models (e.g. \citealp{henderson2023self, rosati2024representation, tamirisa2408tamper, sheshadri2024latent}), existing techniques can consistently be undone within several hundred fine-tuning steps or fewer (e.g. \citealp{qi2024evaluating, huang2025virus, che2025model, fan2025towards, hu2024jogging, sheshadri2024latent, lucki2024adversarial, qian2025layered, deeb2024unlearning}).

\begin{figure}[t]
    \centering
    \includegraphics[width=0.9\linewidth]{figs/fig1b.pdf}
    \caption{\textbf{Training data filtering makes LLMs resistant to adversarial fine-tuning without sacrificing general performance.} Models whose training data has been filtered to remove text related to dual-use biology topics (left) have unaffected general capabilities and (right) have low biothreat proxy capabilities and resist up to 10,000 steps and 300M tokens of adversarial fine-tuning. We further detail results in \Cref{sec:tamper_resistant}.
    }
    \label{fig:fig1}
\end{figure}

% In contrast, open-weight models are typically publicly downloadable for users to run locally and modify. 
% Notable open-weights models include  Llama \citep{grattafiori2024llama, meta2025llama4} and DeepSeek R1 \citep{guo2025deepseek}. 
% Open-weight models are being released at an increasing rate \citep{bhandari2025forecasting}, and their capabilities tend to lag only six to twelve months behind those of closed-weights models \citep{epoch2024openmodelsreport, maslej2024artificial}.
% They offer unique benefits related to transparency, open research, and the decentralization of power \citep{bommasani2024considerations, kapoor2024societal, françois2025differentapproachaisafety}. 
% However, they create unique risks from downstream modifications and unmonitored use \citep{seger2023open, chan2023hazards, huang2024harmful}. 
% However, standard safeguards optimized for closed-weights models do not extend to open-weights models since they do not optimize for tamper-resistance, a critical threat vector enabled by the ability for users to directly modify open-weights models. 


% \stella{Should this paragraph be about finetuning specifically or open weight safety in general? There's non-FT based work we could plug such as ``Suppressing pink elephants with direct principle feedback''} Despite the rising prominence of open-weights models, the science of open-weights model safety is nascent \citep{françois2025differentapproachaisafety}. 
% While closed-weights safety strategies can leverage a provider's full control over weights, prompts, and scaffolding, no such affordances can be relied upon for open-weights models.
% % As a result, open-weights model risk management is challenging, and progress has not been commensurate with the increasing risks and opportunities from open-weights models. 
% % For example, the emerging literature on \textit{AI Safety Cases} \citep{clymer2024safety, buhl2024safety, korbak2025sketch, goemans2024safety, clymer2025example} has yet to propose any strategies for developing strong safety assurances for open-weights systems. 
% The defining challenge for open-weights model safety is that open-weights models can be modified arbitrarily by downstream actors.
% Despite recent work on developing \textit{tamper-resistant} open-weights models (e.g. \citealp{henderson2023self, rosati2024representation, tamirisa2408tamper, sheshadri2024latent}), these techniques can consistently be undone within several hundred fine-tuning steps or fewer (e.g. \citealp{qi2024evaluating, huang2025virus, che2025model, fan2025towards, hu2024jogging, sheshadri2024latent, lucki2024adversarial, qian2025layered, deeb2024unlearning}).
% A key limitation of existing methods for tamper resistance has been their focus on post-training techniques that leverage relatively small amounts of data \citep{jain2023mechanistically, anwar2024foundational}. For example, in this paper, we work with models pretrained on over 1,000 times the number of tokens (550B) used in a typical post-training intervention for tamper-resistance ($\le$150M).

% Despite the rising prominence of open-weights models, the science of open-weights model safety is nascent \citep{françois2025differentapproachaisafety}. 
% The defining challenge is that open-weights models can be modified arbitrarily by downstream actors. 
% Standard safeguards optimized for closed-weights models do not extend to open-weights models since they do not optimize for resistance to harmful tampering, enabled by the ability for users to directly modify the underlying model. 
% As a result, open-weight model risk management is challenging, and progress has not been commensurate with the increasing risks and opportunities from open-weights models. Despite recent work on developing \textit{tamper-resistant} open-weights models (e.g. \citealp{henderson2023self, rosati2024representation, tamirisa2408tamper, sheshadri2024latent}), these techniques can consistently be undone within several hundred fine-tuning steps or fewer (e.g. \citealp{qi2024evaluating, huang2025virus, che2025model, fan2025towards, hu2024jogging, sheshadri2024latent, lucki2024adversarial, qian2025layered, deeb2024unlearning}).

In this paper, we investigate whether open-weight models can be made more tamper-resistant by filtering dual-use content from their pretraining data.
Specifically, we work to make language models robustly ignorant of biothreat proxy knowledge \citep{Li2024TheWB} without degrading unrelated capabilities. 
We hypothesize that if a model is unable to learn unsafe knowledge during pretraining, it will be much more difficult for an attacker to elicit harmful behaviors from it. 
% We distinguish between practical tamper resistance and perfect robustness. 
We make four key contributions:

% We distinguish between practical tamper resistance and perfect robustness. This work does not attempt to address the risks posed by high-resource attackers, such as nation states with significant data and computational resources. Rather, we focus on thwarting low-resource attackers by increasing the cost of tampering. For the purpose of this work, we consider a safeguard to be tamper-resistant if performance on the proxy knowledge is still below the baseline (no safeguards) performance even after training on thousands of proxy documents. 

\begin{enumerate}
    \item \textbf{Knowledge Prevention:} We introduce an efficient multi-stage data filtering pipeline that accounts for less than 1\% of total training FLOPS. We use this filtering approach to successfully prevent biothreat proxy capabilities competitively with existing post-training safeguards. We do not observe degradation to unrelated capabilities
    % . These results suggest that preventing LLMs from learning dual-use technical knowledge through data curation is both tractable and effective 
    (\Cref{sec:filtering}).
    \item \textbf{Tamper-Resistance:} We show that training data filtering\footnote{Filtering \& Evaluation code is available at \href{https://github.com/EleutherAI/deep-ignorance}{this link}.} can achieve \textit{state-of-the-art} tamper resistance for up to 10,000 steps and 300M tokens of adversarial fine-tuning on biothreat-related text, improving by more than an order of magnitude over post-training baselines (\Cref{sec:tamper_resistant}).
    \item \textbf{Defense-in-Depth:} We demonstrate that data filtering cannot prevent LLMs from leveraging harmful knowledge provided in-context, but that Circuit-Breaking-based techniques \citep{zou2024improving} offer complementary defenses. However, we show that none of the defenses we test are resistant to staged attacks that combine fine-tuning and in-context retrieval (\Cref{sec:defense_in_depth}).
    \item \textbf{Model Suite:} We release a suite of 6.9B parameter language models trained with combinations of data filtering and post-training safeguards. 
    These models will allow researchers to study the causal impact that removing a subset of training data has on model mechanisms and behavior.\footnote{Data \& models are available at \href{https://huggingface.co/collections/EleutherAI/deep-ignorance-685441040d024a0fee593d68}{this link}. We overview released models in \Cref{tab:released_models}.}
\end{enumerate}

In addition to these core contributions, we also study challenges with evaluating LLM safeguards using multiple-choice questions, introduce a new hybrid unlearning algorithm that combines prior state-of-the-art defenses \citep{zou2024improving, sheshadri2024latent}, and present negative results involving our challenges with synthetic document training \citep{anthropic2025modifying} at scale. 
Finally, while this work does not focus on machine unlearning, our models offer a test-bed for studying machine unlearning algorithms.
Together, we hope that these findings can contribute to emerging science and standards for effective open-weight model risk management. 

% \begin{figure}[t!]
%     \centering
%     \includegraphics[width=\linewidth]{figs/fig1.pdf}
%     \caption{Pretraining data filtering makes LLMs resistant to unwanted fine-tuning. (a) Whereas most research on suppressing harmful model capabilities focuses on post-training interventions, we filter during pretraining and fine-tuning (600B tokens) instead of post-training ($\approx$200M tokens). (b) filtering biothreat-related conetnt from trainign data  can make models effectively resist harmful fine-tuning for up to 300M tokens on biothreat-related papers.}
%     \label{fig:fig1}
% \end{figure}

% List of main experiments
% \kyle{
%     This list captures all the experiments we want to discuss in the main paper. 
%     \begin{enumerate}
%         \item Compare filtering setups (4 combinations) vs circuit breaking and baseline. No tampering.
%         \item Compare filtering setups (4 combinations) vs circuit breaking and baseline. This time, we measure performance after applying varying amounts of adversarial fine-tuning on proxy knowledge. 
%         \item Compare filtering setups (4 combinations) vs circuit breaking and baseline. This time, we measure performance after applying varying amounts of fine-tuning on benign/general text. We expect this to undo circuit breakers but keep filtering performance flat. 
%         \item Compare all approaches on a new synthetic dataset with open and closed-book options. 
%     \end{enumerate}
% }
